\vfill\pagebreak
\section{A first look at \texttt{match}}
\label{sec:control:match}

A chain of \token{else if} that compares one value with several possibilities
reads badly: the value is repeated in every condition, and the eye has to check
that it really is the same value each time. \token{match} states the value once
and lists the possibilities below it.

\begin{lstlisting}[style=coloredverbatim, caption=Choosing on the value of a number, label=lst:control:match]
use std::io;

fn main() {
  let day = read!i32(ask-> "Day of the week (1 to 7): ");
  match day {
    6 | 7 => {
      println("Weekend.");
    }
    1 ... 5 => {
      println("Weekday.");
    }
    _ => {
      println("There is no day ", day, ".");
    }
  }
}
\end{lstlisting}
\vspace{-10pt}%
\begin{lstlisting}[style=bashVerb, escapechar=@+]
@+\textcolor{teal!80}{alice@dev:\mtilde}@+$ ./main
Day of the week (1 to 7): 6
Weekend.
@+\textcolor{teal!80}{alice@dev:\mtilde}@+$ ./main
Day of the week (1 to 7): 9
There is no day 9.
\end{lstlisting}

Each line of the form \token{pattern => \{ ... \}} is an \textit{arm}. The arms
are tried from top to bottom, as the conditions of an \token{else if} chain are,
and the first one whose pattern fits the value runs its block. The others are
skipped. The arms are not separated by commas.

The patterns of Listing~\ref{lst:control:match} show the three forms that work on
numbers:

\begin{itemize}
\item a value, such as \token{7}, which fits only that value; several values
  separated by \token{|} fit any one of them, so \token{6 | 7} is the weekend;
\item a range, such as \token{1 ... 5}, which fits every number the range
  contains, with the same two operators as in Section~\ref{sec:control:for};
\item the wildcard \token{\_}, which fits anything, and is the \token{else} of a
  \token{match}.
\end{itemize}

Since \token{\_} fits anything, an arm written after it could never run. The
compiler rejects such an arm with ``matcher expression is never evaluated''. The
catch-all arm always comes last.

\subsection{Choosing a value}

A \token{match}, like an \token{if}, is an expression, and each arm can produce
a value instead of running instructions. When an arm is a single value, its
braces can go.

\begin{lstlisting}[style=coloredverbatim, caption=A \token{match} used as a value, label=lst:control:match_value]
use std::io;

fn main() {
  let day = read!i32(ask-> "Day of the week (1 to 7): ");
  let name = match day {
    1 => "Monday"
    2 => "Tuesday"
    3 => "Wednesday"
    4 => "Thursday"
    5 => "Friday"
    6 => "Saturday"
    7 => "Sunday"
    _ => "no day at all"
  };
  println("Day ", day, " is ", name, ".");
}
\end{lstlisting}

The rule is the one of Section~\ref{sec:control:if}: every possible run must
leave \token{name} with a value. Without the last arm, a day of 9 would match
nothing, so a \token{match} used as a value must end with an arm that catches
everything. Leaving it out is an error: ``the pattern matching has no default
case''.

\subsection{Naming the value: bindings and guards}

An arm can give a name to the value it matches, so that its block can use it.
The name declares a variable, which holds the matched value and exists only in
that arm. It is written in one of two ways.

\begin{itemize}
\item \token{name = pattern} names the value when the pattern after \token{=}
  fits it. An arm whose pattern is \token{s = 1 ... 9} fits the numbers from 1
  to 9, and inside its block \token{s} is the one that matched.
\item A name alone is a pattern of its own. Like \token{\_}, it fits any value,
  but unlike \token{\_} it gives the arm's block that value to work with.
\end{itemize}

An arm can also test the value further. After the pattern, \token{if} followed
by a condition adds a \textit{guard}, and the arm fits only when the guard holds
too. The condition can use the name that the pattern declared.
Listing~\ref{lst:control:match_guard} uses a binding on line~10, and a name
with a guard on line~13.

\begin{lstlisting}[style=coloredverbatim, caption=Arms with a binding and a guard, label=lst:control:match_guard]
use std::io;

fn main() {
  let limit = read!i32(ask-> "Speed limit: ");
  let speed = read!i32(ask-> "Your speed: ");
  match speed {
    0 => {
      println("You are not moving.");
    }
    s = 1 ... 9 => {
      println("Only ", s, "? You could walk faster.");
    }
    s if s > limit => {
      println("Slow down, you are ", s - limit, " over the limit.");
    }
    _ => {
      println("All good.");
    }
  }
}
\end{lstlisting}
\vspace{-10pt}%
\begin{lstlisting}[style=bashVerb, escapechar=@+]
@+\textcolor{teal!80}{alice@dev:\mtilde}@+$ ./main
Speed limit: 50
Your speed: 62
Slow down, you are 12 over the limit.
@+\textcolor{teal!80}{alice@dev:\mtilde}@+$ ./main
Speed limit: 50
Your speed: 5
Only 5? You could walk faster.
\end{lstlisting}

A name in a pattern always declares a new variable; it never reads an existing
one, and that holds on either side of \token{=}. So an arm
\token{s = limit => \{ ... \}} does not mean ``when the speed equals the
limit''. It declares two variables, \token{s} and \token{limit}, both holding
the speed, and since the program already has a \token{limit}, the compiler stops
with ``declaration of limit shadows another declaration''. A comparison with a
variable goes in a guard: the arm that means ``equal to the limit'' is
\token{s if s == limit => \{ ... \}}.

This chapter matches integers only. \token{match} is far more capable than
that -- it can take values apart, as later chapters show -- and
Section~\ref{sec:flow:pattern_matching} lists everything it can do.
